\chapter{Introduction}

\section{Motivation}

The rapid evolution of Artificial Intelligence (AI) has fundamentally change the global socio-economic landscape. To adapt to the new era, self\-studying and learning method are becoming more and more important. Traditional education philosophy, such as "practice makes perfect," has no longer been the absolute true; when AI can automate execution, the value of learning moves from procedural fluency and labour expertise to conceptual mastery and strategic analysis. However, current education system fail to keep pace with this evolution:

\begin{itemize}
    \item \textbf{Stagnancy of Content:} Legacy Learning Management Systems (LMS) treat knowledge as static artifacts (textbooks, slides). These materials fail to adapt to the varying cognitive baselines of individual students.
    \item \textbf{The Guidance Deficit:} Without a dynamic roadmap, students navigate a "sea of information" where external resources are either too granular (leading to information overload) or too condensed (creating knowledge gaps), with little profound ground truth.
    \item \textbf{The Lack of Interaction:} Current systems lack a feedback loop. Textbooks exceed the average human attention span, while slides provide too condensed context insufficient for novices, leaving no room for real-time, human-like inquiry or mentorship.
    \item \textbf{Outdated Knowledge:} Educational materials and methods cannot catch up with the change of technology. Internet are an invaluable resource, however information is chaotic, unstructured and content are confusing to many new learners.
\end{itemize}

While Large Language Models (LLMs) offer a glimpse into automated resource synthesis and knowledge base construction, significant technical bottleneck prevent their reliable application in complex domains:

\begin{itemize}
    \item \textbf{Semantic Fragmentation in RAG:} Traditional Retrieval-Augmented Generation (RAG) frameworks rely on discrete data chunking. This fragmentation disrupts the global structural integrity of a subject, leading to "contextual myopia" where the model captures isolated facts but fails to comprehend the overarching hierarchy of the domain.
    \item \textbf{Contextual Saturation vs. Precision:} Expanding the context window of LLMs often introduces diminishing returns. The injection of unrefined data frequently leads to "contextual overhead" and "noise," where the model struggles to distinguish fundamental principles from peripheral details.
    \item \textbf{Relational Limitations of Vector Space:} Relying exclusively on vector embeddings (typically $\approx 1024$ dimensions) is insufficient for mapping the non-linear, multi-faceted dependencies inherent in academic taxonomies. Similarity-based retrieval cannot substitute for the deterministic logic required to understand prerequisite relationships.
    \item \textbf{The Reasoning Gap in Pedagogy:} Instruction is a sophisticated, non-linear reasoning task. Current AI systems are largely reactive—proficient in text generation but deficient in long-term goal tracking and proactive guidance. This results in a susceptibility to hallucinations and an inability to execute the "evaluate-and-guide" workflow essential for effective pedagogy.
\end{itemize}

\section{Goals}

This project aims to bridge the gap between passive content and active learning by developing a personalized AI-driven educational ecosystem. The specific objectives include:

\begin{itemize}
    \item \textbf{Philosophy of Co-living with AI:} Transforming AI from a simple tool into a pervasive personal assistant that evolves alongside the learner.
    \item \textbf{Cognitive State Tracking:} Establishing a system to record student progress, identify long-term goals, and objectively evaluate theoretical comprehension through interaction.
    \item \textbf{Adaptive Learning Trajectories:} Customizing the learning pace and curriculum dynamically, ensuring that the difficulty level remains within the student's "Zone of Proximal Development."
\end{itemize}

\section{Scope}

To fulfill these goals, the research focuses on the following technical cores:

\begin{itemize}
    \item \textbf{Knowledge Modeling:} Developing a deterministic "Backbone Graph" that maps semantic connections and context grouping, serving as a structural pattern for LLM reasoning. 
    \item \textbf{Context Engineering:} Engineering methods to facilitate reasoning across vast, high-density knowledge bases without losing global coherence.
    \item \textbf{Agent Workflow Orchestration:} Designing multi-agent systems to handle complex, multi-step pedagogical reasoning tasks.
    \item \textbf{Systemic Architecture:} Implementing clean design patterns optimized for scalable, AI-native educational services.
    \item \textbf{User Experience (UX):} Creating a dual-interface system featuring real-time agent streaming for interaction and a comprehensive GUI for administrative and roadmap visualization.
\end{itemize}


\section{Thesis structure}
There are 7 chapters in this project's report:

\begin{itemize}
    \item \textbf{Chapter 1} introduces the problems, proposes the system's goals, scopes and thesis structure
    \item \textbf{Chapter 2} provides fundamental knowledge for project understand and deep literature review of relevant prior researches covering topic, as well as preview of open-source application.
    \item \textbf{Chapter 3} Requirement Specifications
    \item \textbf{Chapter 4} propose methodologies: we discuss what inferred from previous work and what ways to solve the specific usecase of ours
    \item \textbf{Chapter 5} details the system including: System architecture and its components, settings in each components, and finally, their implementation explanation that fully illustrate how do the system put proposed methodology to practice 
    \item \textbf{Chapter 6} analyzes the system outcomes from components to final output as a whole
    \item \textbf{Chapter 7} is the conclusion to the result: what the system achieved, its limitation and future improvements for the final Capstone Project
\end{itemize}